\documentclass[11pt,a4paper]{article}

\usepackage[margin=2.4cm]{geometry}
\usepackage{amsmath,amssymb}
\usepackage{booktabs}
\usepackage{longtable}
\usepackage{xcolor}
\usepackage{listings}
\usepackage[colorlinks=true,linkcolor=blue!50!black,urlcolor=blue!50!black,
            citecolor=blue!50!black]{hyperref}
\usepackage{titlesec}
\usepackage{enumitem}

\definecolor{codebg}{gray}{0.96}
\definecolor{keyw}{rgb}{0.15,0.35,0.6}
\lstset{
  basicstyle=\ttfamily\footnotesize,
  backgroundcolor=\color{codebg},
  frame=single, framesep=4pt, rulecolor=\color{gray!40},
  breaklines=true, breakatwhitespace=false,
  columns=fullflexible, keepspaces=true,
  showstringspaces=false,
  xleftmargin=0pt, xrightmargin=0pt,
  literate={μ}{{$\mu$}}1 {ω}{{$\omega$}}1 {→}{{$\rightarrow$}}1 {×}{{$\times$}}1
}

\newcommand{\code}[1]{\texttt{\small #1}}
\newcommand{\wpt}{\omega_p t}
\newcommand{\vth}{v_{\rm th}^{e}}
\newcommand{\vthsq}{v_{{\rm th},e}^{2}}
\newcommand{\Npc}{N_{\rm pc}}

\title{\bfseries Reproducing \textit{``No cosmological constraints on dark photon
dark matter from resonant conversion''} (arXiv:2510.13956) with \textsc{Entity}\\[4pt]
\large Implementation on Lonestar6 and comparison with the published runs}
\author{}
\date{\today}

\begin{document}
\maketitle

\begin{abstract}
\noindent
This note documents a reproduction of the particle-in-cell (PIC) results of Hook,
Huang \& Shalaby (arXiv:2510.13956v1) using the \textsc{Entity} PIC framework in
place of the authors' \textsc{SHARP} code.

Part~\ref{part:guide} is an operating manual: the repository layout, a quickstart, a
reference for the input deck, job scripts and queue selection, the output format and
analysis pipeline, a procedure for sizing a new run, and a consolidated list of the
rules and traps that govern all of it. It is reference-ordered and is the only part
needed in order to run something.
Part~\ref{part:impl} covers the implementation: how a problem generator attaches to
the framework, built up from a minimal working example; the generator written for this
work; a deterministic (``quiet'') particle loader; the four build
configurations required on TACC Lonestar6; and a measured performance characterisation
across CPU, single-GPU and multi-GPU execution. Part~\ref{part:runs} documents every
production run, the exact settings used, and a point-by-point comparison with the
parameters and results stated in the paper.

Both resonance cases of the paper's Fig.~2 are reproduced. The Landau--Zener case is
not: we measure an electron heating of ${\sim}10^3$ where the paper reports a factor
below $2$, and four candidate explanations have been tested and excluded. Since the
resonance cases agree, the disagreement is specific to the swept drive. All
performance and physics numbers quoted are measured on Lonestar6 unless explicitly
labelled as an estimate; \S\ref{sec:corrections} lists the figures from earlier drafts
that are superseded, several of which were artefacts of a particle loader that
satisfied the paper's stated initial condition but not its subsequent behaviour.
\end{abstract}

\tableofcontents
\bigskip\hrule\bigskip

\section{Scope and conventions}

The paper studies the resonant conversion of dark-photon dark matter (DPDM) into a
plasma excitation in the early-universe electron--ion plasma, and shows that
\emph{nonlinear} plasma dynamics saturate the conversion far below the linear-theory
expectation, weakening the resulting cosmological limits. The evidence is a set of
1D--1V electrostatic PIC simulations.

Stripped of the cosmology, the numerical problem is small and sharply defined:

\begin{itemize}[nosep]
  \item 1D electrostatic, periodic box, $N_x = 1000$, $L = 40\,c/\omega_p$,
        $\Delta x = 0.04\,c/\omega_p \simeq \lambda_D$;
  \item electrons and ions with the \emph{real} mass ratio $m_i/m_e = 1836$, both
        evolved;
  \item $k_B T_e = k_B T_i = 10^{-3} m_e c^2$;
  \item the dark photon enters as a \emph{spatially uniform} ($k=0$), time-oscillating
        external electric field $E_{\rm ext}(t) = A_0\cos(\omega t)$ acting on the
        particle Lorentz force \emph{only} --- it is never a source term in Maxwell's
        equations. All feedback is through the deposited current.
\end{itemize}

\paragraph{Code units.} Throughout, $\omega_{pe} = 1$, $c = 1$, $n_e = n_0$,
$m_e = 1$. Energies are in $n_e m_e c^2$, times in $\omega_p^{-1}$, wavenumbers in
$\omega_p/c$. The drive amplitude in code units equals the paper's:
$A_0 = v_q^D/c = (v_q^D/\vth)\sqrt{10^{-3}}$.

\paragraph{Measured versus estimated.} Every performance and physics number here is
measured on Lonestar6 unless explicitly labelled an estimate.
Section~\ref{sec:corrections} lists the few figures from earlier drafts that are
superseded, so that they are recognisable if they turn up in older scripts or notes.

\clearpage
\part{Working guide}
\label{part:guide}

\noindent\emph{This part is reference-ordered and self-contained: it is what to read in
order to run something. Part~\ref{part:impl} explains why the code is built the way it
is, and Part~\ref{part:runs} records what the production runs produced. Neither is
needed to operate the code.}

\section{Repository map}
\label{sec:repomap}

Everything specific to this work lives in two directories of the \textsc{Entity} tree.

\begin{center}\small
\begin{tabular}{@{}lp{9.6cm}@{}}
\toprule
Path (under repository root) & Contents \\
\midrule
\multicolumn{2}{@{}l}{\textit{\code{pgens/dpdm/} --- problem generator, decks, job scripts}}\\
\code{pgen.hpp}            & The problem generator: external drive, quiet loader, custom
                             diagnostic. The only C++ in this work (\S\ref{sec:pgen}). \\
\code{resonance.toml}      & Reference deck, fixed-frequency family. Start here. \\
\code{landau\_zener.toml}  & Reference deck, swept-frequency family. \\
\code{run\_paper\_fast.sh} & Production: $v_q^D/\vth = 3\times10^{-2}$, $\Npc=2\times10^2$, 1 CPU node. \\
\code{run\_paper\_slow.sh} & Production: $v_q^D/\vth = 1\times10^{-3}$, $\Npc=2\times10^5$, 1 A100. \\
\code{run\_paper\_slow\_resume.sh} & The same run restarted from checkpoint (\S\ref{sec:restart}). \\
\code{run\_lz1p0.sh}, \code{run\_lz0p1.sh} & Landau--Zener production at $v_q^D/\vth = 1.0$ and $0.1$. \\
\code{run\_gate\_loader.sh} & The loader gate: three undriven runs, one per
                             \code{loading}, against the paper's Fig.~A00
                             (\S\ref{sec:quietval}). \\
\code{run\_paper\_fast\_shear.sh}, \code{run\_paper\_slow\_shear.sh},
\code{run\_lz1p0\_shear.sh} & The three production re-runs with the corrected
                             loader. Identical decks; new output roots. \\
\code{run\_lz1p0\_naive.sh} & Landau--Zener with \code{sweep\_phase = "naive"}. \\
\code{run\_lz\_heating.sh} & Undriven numerical-heating gate at Landau--Zener duration. \\
\code{validate\_quiet.sh}, \code{check\_quiet.py} & Quiet-start validation gates (\S\ref{sec:quietval}). \\
\code{validate\_mpi.sh}    & Confirms results are independent of MPI domain count. \\
\code{bench\_slow.sh}, \code{bench\_gpu.sh}, \code{bench\_multigpu.sh} & Throughput benchmarks: CPU/MPI, 1 GPU, 3 GPU. \\
\code{diag\_slow.sh}, \code{diag\_growth.sh}, \code{affinity\_test.sh} & Diagnostics for
                             performance and thread-placement questions. \\
\midrule
\multicolumn{2}{@{}l}{\textit{\code{pgens/dpdm/paper/} --- analysis pipeline for the paper-faithful runs}}\\
\code{extract\_paper.py}   & Reduces \code{fast}, \code{slow}, \code{lz1p0}, \code{heat}
                             and the quiet/random pair to \code{.npz}
                             (\S\ref{sec:analysis}). \\
\code{PLOTS.md}            & The six plots those runs support, and what each claims. \\
\midrule
\multicolumn{2}{@{}l}{\textit{\code{pgens/dpdm/run34/} --- the resonance-ladder survey and its analysis pipeline}}\\
\code{run34.sh}            & Eight-run ladder in $v_q^D/\vth$ plus two convergence runs. \\
\code{extract\_run34.py}   & Reduces ADIOS2 output to \code{.npz} (\S\ref{sec:analysis}). \\
\code{PLOTS.md}            & Which four plots to make and what each one claims. \\
\code{run34\_plots.ipynb}  & The plots, from the \code{.npz} files. \\
\midrule
\multicolumn{2}{@{}l}{\textit{\code{dpdm\_doc/} --- this document}}\\
\code{main.tex}, \code{Makefile} & Source and build (\code{make} $\to$ \code{main.pdf}). \\
\bottomrule
\end{tabular}
\end{center}

\noindent Build directories (\code{build-dpdm*/}) are deliberately not version
controlled; see \S\ref{sec:builds}. Run output goes to \code{\$SCRATCH}, which is
\textbf{purged on a rolling basis} --- copy anything you intend to keep to
\code{\$WORK}.

\section{Quickstart}
\label{sec:quickstart}

End to end, on a login node unless stated otherwise:

\begin{enumerate}
  \item \textbf{Build once.} Follow \S\ref{sec:builds}. For a first run the CPU
        OpenMP configuration (\code{build-dpdm}) is enough; the CUDA build
        (\code{build-dpdm-cuda}) is what production uses.
  \item \textbf{Verify the build} --- configure options fail silently here:
        \begin{lstlisting}
grep -E "^Kokkos_ENABLE_(OPENMP|SERIAL|CUDA):BOOL" build-*/CMakeCache.txt
        \end{lstlisting}
  \item \textbf{Copy a job script} from \code{pgens/dpdm/} that matches the shape of
        what you want --- they are self-contained and embed their own deck --- and edit
        the deck inside it (\S\ref{sec:deck}).
  \item \textbf{Size the run}: particle count, GPU count and walltime follow a fixed
        procedure (\S\ref{sec:sizing}). The GPU count in particular cannot be changed
        after the first submission if the run will ever restart.
  \item \textbf{Submit}: \code{sbatch run\_my\_case.sh}. Jobs under 2\,h belong in
        \code{development} (\S\ref{sec:queues}).
  \item \textbf{Analyse}: the stats CSV answers most questions directly; field and
        particle dumps go through \code{extract\_run34.py}
        (\S\ref{sec:analysis}).
\end{enumerate}

\section{The input deck}
\label{sec:deck}

\subsection{The \code{[setup]} block --- the DPDM interface}

These nine keys are the whole problem-specific interface; everything else in the deck
is stock \textsc{Entity}. Defaults are those in \code{pgen.hpp}.

\begin{center}\small
\begin{tabular}{@{}llp{7.4cm}@{}}
\toprule
Key & Default & Meaning \\
\midrule
\code{drive}        & \code{"resonance"} & \code{"resonance"} (fixed $\omega$) or
                                           \code{"landau\_zener"} (swept). \\
\code{loading}      & \code{"random"}    & \code{"quiet"} is the paper's IC:
                                           lattice positions, sampled velocities,
                                           relaxing to the Poisson floor.
                                           \code{"quiet\_lattice"} is the legacy
                                           permanently-ordered loader, kept only to
                                           reproduce superseded runs --- it does
                                           \emph{not} reproduce the paper
                                           (\S\ref{sec:quietrelax}).
                                           \code{"random"} is the stock injector and
                                           has $\epsilon_{\rm noise}(0)\neq0$.
                                           \textbf{Set this explicitly.} \\
\code{one\_v}       & \code{true}        & Sets $u_{x2}=u_{x3}=0$, matching
                                           \textsc{SHARP} and the paper's 1V thermal
                                           speed. Leave on for paper comparison. \\
\code{A0}           & \code{0.0}         & Drive amplitude in code units,
                                           $A_0 = v_q^D/c = (v_q^D/\vth)\sqrt{10^{-3}}$. \\
\code{omega}        & \code{1.0}         & Drive frequency, \code{drive="resonance"} only. \\
\code{sweep\_phase} & \code{"integrated"} & \code{"integrated"} gives instantaneous
                                           frequency $\omega_0+\dot\omega t$ and is
                                           correct; \code{"naive"} reproduces the
                                           literal reading of the paper's notation
                                           and sweeps at twice the rate
                                           (\S\ref{sec:sweep}). \code{landau\_zener}
                                           only. \\
\code{omega0}       & \code{0.0}         & Initial frequency, \code{drive="landau\_zener"} only. \\
\code{domega\_dt}   & \code{0.0}         & Sweep rate, \code{drive="landau\_zener"} only. \\
\code{densities}    & \code{[2.0]}       & One entry, per species \emph{pair}.
                                           \code{[2.0]} with \code{skindepth0 = 1}
                                           gives $n_e = n_0$, hence $\omega_{pe}=1$. \\
\code{temperatures} & ---                & One entry per species, in $m_0c^2$. The
                                           injector divides by species mass, so
                                           \code{[1.0e-3, 1.0e-3]} is
                                           $10^{-3}m_ec^2$ for both. \\
\bottomrule
\end{tabular}
\end{center}

\noindent The two drive families are mutually exclusive in their parameters:
\code{omega} is ignored when sweeping, and \code{omega0}/\code{domega\_dt} are ignored
when not. An unrecognised \code{drive} or \code{loading} aborts at startup with a
readable message rather than running the wrong physics.

\subsection{Settings outside \code{[setup]} that are not optional}
\label{sec:deckmust}

\begin{itemize}[nosep]
  \item \code{particles.ppc0} must be written as a \textbf{float}.
        \code{ppc0 = 2000} aborts with
        \code{toml::value::as\_floating(): bad\_cast}; write \code{2000.0}.
  \item \code{maxnpart} per species must exceed
        $\code{ppc0}\times\code{densities[0]}/2\times N_x$, with headroom.
  \item A \code{[diagnostics]} block is required in practice.
        Without it, \code{<name>.log} is a kernel trace at ${\sim}4$\,kB
        \emph{per step} and \code{<name>.out} is a progress bar redrawn every step,
        both written synchronously to the parallel filesystem. Use
        \code{log\_level = "WARNING"}, \code{interval = 1000},
        \code{colored\_stdout = false}.
  \item \code{[output.spectra]} defaults to \emph{on}; set \code{enable = false}.
  \item \code{[output.stats] interval\_time} must be $\ge 0.5$, and preferably
        $\ge 1.0$. Per-species \code{T00\_1}/\code{T00\_2} force a moment reduction
        over every particle; at $0.1$ this doubles wall-clock.
  \item Stats quantity names are literal: \code{"E\^{}2"}, \code{"T00\_1"},
        \code{"T00\_2"}, \code{"J.E"}. \code{"JdotE"} aborts. \code{"T00"} alone gives
        the species sum.
  \item \code{algorithms.current\_filters = 0}. The saturation physics lives at high
        $k$; filtering removes exactly the modes of interest.
\end{itemize}

\subsection{A complete deck}

The fast resonance case in full. Every other run in this work is this file with
\code{ppc0}, \code{A0}, \code{runtime} and the \code{[setup]} drive keys changed.

\begin{lstlisting}
[simulation]
  name = "fast"
  engine = "srpic"
  runtime = 10000.0            # omega_p t
[grid]
  resolution = [1000]          # Nx
  extent = [[0.0, 40.0]]       # L = 40 c/omega_p, so dx = 0.04 ~ lambda_D
  [grid.metric]
    metric = "minkowski"
  [grid.boundaries]
    fields = [["PERIODIC"]]
    particles = [["PERIODIC"]]
[scales]
  skindepth0 = 1.0             # with densities=[2.0] this sets omega_pe = 1
  larmor0 = 1.0                # unused: no magnetic field
[algorithms]
  current_filters = 0          # KEEP AT 0 -- see above
  [algorithms.timestep]
    CFL = 0.5                  # dt = 0.02 / omega_p
[particles]
  ppc0 = 200.0                 # N_pc per cell per species; Np = N_pc*Nx = 2e5
  [[particles.species]]
    label = "e-"
    mass = 1.0
    charge = -1.0
    maxnpart = 2.5e5
    pusher = "Vay"
  [[particles.species]]
    label = "i+"
    mass = 1836.0              # real mass ratio; ions are evolved
    charge = 1.0
    maxnpart = 2.5e5
    pusher = "Vay"
[setup]
  densities = [2.0]
  temperatures = [1.0e-3, 1.0e-3]
  drive = "resonance"
  loading = "quiet"            # NOT the default -- must be stated.
                               # "quiet" = lattice positions + sampled velocities.
                               # "quiet_lattice" is the legacy loader; it never
                               # decoheres and does not reproduce the paper.
  one_v = true
  A0 = 9.48683e-4              # v_q^D/vth_e = 3e-2
  omega = 1.0
[diagnostics]
  interval = 1000
  log_level = "WARNING"
  colored_stdout = false
[output]
  interval_time = 10.0         # field/particle dumps; 1000 of them here
  [output.fields]
    quantities = ["E", "N_1", "N_2", "Rho", "J", "T00_1", "T00_2"]
  [output.particles]
    species = [1, 2]
    stride = 20
  [output.stats]
    enable = true
    interval_time = 0.5
    quantities = ["E^2", "T00_1", "T00_2", "J.E"]
    custom = ["e_ext"]         # the drive itself, for phase checks
  [output.spectra]
    enable = false
[checkpoint]
  interval_time = 2500.0
  keep = 2
\end{lstlisting}

\noindent For the swept family, replace the drive keys:
\begin{lstlisting}
  drive = "landau_zener"
  A0 = 3.16228e-3              # v_q^D/vth_e = 0.1
  omega0 = 0.8                 # start below resonance
  domega_dt = 2.0e-5           # 0.1 omega_p over omega_p t = 5000
  sweep_phase = "integrated"   # crossing at omega_p t = 1e4;
                               # "naive" doubles the rate -> crossing at 5000
\end{lstlisting}

\section{Job scripts and the run loop}
\label{sec:jobs}

\subsection{Template}

Job scripts in \code{pgens/dpdm/} are self-contained: they carry their own deck as a
heredoc, so a script plus its log is a complete record of what was run. The skeleton:

\begin{lstlisting}
#!/bin/bash
#SBATCH -J dpdm_fast
#SBATCH -o sbatch_logs/out/fast.o%j
#SBATCH -e sbatch_logs/err/fast.e%j
#SBATCH -p development          # see queue selection below
#SBATCH -N 1
#SBATCH -n 1
#SBATCH -t 01:30:00
#SBATCH -A PHY23028
#SBATCH --mail-type=all
#SBATCH --mail-user=<you>@utexas.edu

module load gcc/13.2.0          # REQUIRED at run time, not only at build time

export OMP_NUM_THREADS=112      # CPU builds only; omit for CUDA
export OMP_PROC_BIND=close
export OMP_PLACES=cores

ROOT=$SCRATCH/dpdm/paper/fast
EXE=<repo>/build-dpdm/src/entity.xc
mkdir -p $ROOT/fast             # pre-create <simulation.name>/ -- see below

cat > $ROOT/fast.toml <<EOF
... deck ...
EOF

cd $ROOT && $EXE -input fast.toml > fast.stdout 2>&1
echo "exit=$? $(date -Is)"
\end{lstlisting}

Three things in that skeleton are load-bearing:

\begin{itemize}[nosep]
  \item \code{module load gcc/13.2.0} inside the job. The binary links GCC~13's
        \code{libstdc++}; the default environment is gcc/9.4.0 and the executable dies
        with \code{GLIBCXX\_3.4.31 not found} \emph{before} reading its input.
  \item \code{mkdir -p \$ROOT/<simulation.name>}. Under MPI every rank calls
        \code{mkdir} on that directory and the losers abort with
        \code{filesystem error: cannot create directory: File exists}, exit 255, within
        seconds. A job that "finished" in four seconds is this, not speed.
  \item \code{sbatch\_logs/out/} and \code{sbatch\_logs/err/} must exist before
        submission or the job is rejected.
\end{itemize}

MPI runs replace the last line with \code{ibrun} and set \code{-n} to the rank count;
single-GPU runs need no launcher and no \code{OMP\_*} variables.

\subsection{Choosing a queue}
\label{sec:queues}

\begin{center}\small
\begin{tabular}{@{}lrrrrr@{}}
\toprule
Partition & MinNode & MaxNode & MaxWall & MaxNodePU & MaxJobsPU \\
\midrule
\code{development}    & 1 & 8   & 02:00:00   & 8  & 1  \\
\code{normal}         & 1 & 64  & 2-00:00:00 & 64 & 20 \\
\code{gpu-a100}       & 1 & 8   & 2-00:00:00 & 12 & 8  \\
\code{gpu-a100-dev}   & 1 & 2   & 02:00:00   & 2  & 1  \\
\code{gpu-a100-small} & 1 & 1   & 2-00:00:00 & 3  & 3  \\
\bottomrule
\end{tabular}
\end{center}

\begin{itemize}[nosep]
  \item \textbf{Under 2\,h and $\le 8$ nodes: use \code{development}.} Measured
        turnaround there was 9\,s and 70\,s for jobs that never started at all in
        \code{normal}. \code{MaxJobsPU = 1} only means short jobs run back-to-back,
        which still beats \code{normal} priority.
  \item \textbf{Single-GPU production: \code{gpu-a100-small}.} One A100 per node,
        billing weight 32, versus three per node at weight 128 on \code{gpu-a100}.
  \item \textbf{Prefer fewer nodes than the scaling curve suggests.} On a busy machine
        queue wait swamps compute: multi-node GPU jobs requesting 2 and 4 nodes sat
        unscheduled for 30\,h and were cancelled, while single-node jobs started within
        minutes. A 12-GPU configuration with a 6\,h runtime is worthless if it waits
        30\,h for contiguous nodes.
\end{itemize}

\subsection{Checkpoints and restarting}
\label{sec:restart}

Checkpoints are controlled by \code{[checkpoint] interval\_time} with \code{keep}
retained, and written to \code{<name>.ckpt/} as
\code{step-NNNNNNNN.bp} plus a \code{meta-NNNNNNNN.toml}. A run resumes with

\begin{lstlisting}
cd $ROOT && $EXE -input fast.toml -resume > fast.resume.stdout 2>&1
\end{lstlisting}

\noindent (aliases: \code{-continue}, \code{-restart}, \code{-checkpoint}).
\textsc{Entity} selects the latest checkpoint automatically. Increase \code{runtime}
in the deck if the original value has already been reached.

\textbf{The domain count is locked by the checkpoint format.} An $N$-domain checkpoint
can only be resumed by $N$ domains: a single-GPU checkpoint cannot be resumed by a
3-GPU binary. Decide the rank/GPU count before the \emph{first} submission of any run
that might need to restart --- there is no way to change it later without starting
over.

\section{Output and analysis}
\label{sec:analysis}

\subsection{What a run writes}

Output lands in \code{<rundir>/<simulation.name>/}:

\begin{center}\small
\begin{tabular}{@{}lp{9.4cm}@{}}
\toprule
File & Contents \\
\midrule
\code{<name>\_stats.csv} & Scalar time series, one row per
   \code{[output.stats] interval\_time}. Answers most questions on its own. \\
\code{fields/*.bp}       & ADIOS2 BP5 field dumps at \code{[output] interval\_time}. \\
\code{particles/*.bp}    & Particle dumps, decimated by \code{[output.particles] stride}. \\
\code{<name>.log}, \code{<name>.out}, \code{<name>.info}, \code{<name>.err} & Logs; see
   \S\ref{sec:deckmust} before the first production run. \\
\code{<name>.ckpt/}      & Checkpoints (sibling of \code{<name>/}, not inside it). \\
\bottomrule
\end{tabular}
\end{center}

The stats CSV columns, in order, are
\code{step}, \code{time}, \code{E1\^{}2}, \code{E2\^{}2}, \code{E3\^{}2},
\code{T00\_1}, \code{T00\_2}, \code{J.E}, and any \code{custom} entries
(here \code{e\_ext}). In 1D electrostatic runs \code{E2\^{}2} and \code{E3\^{}2} are
zero by construction.

\subsection{Deriving the physical quantities}

The conversions below are the ones the plots need, and two of them are easy to get
wrong in a way that changes conclusions by a factor of two or three.

\begin{itemize}[nosep]
  \item $\epsilon_E = \tfrac12\langle E_x^2\rangle$ from \code{E1\^{}2}.
  \item $\Delta KE_e = T00\_1(t) - T00\_1(0)$, and likewise for the ions. Taking the
        difference is what makes this safe under \textsc{Entity}'s 3V convention: with
        no transverse forces the perpendicular energy is constant and cancels.
  \item \textbf{The linear-theory line is
        $\epsilon_E + \Delta KE_e = C A_0^2 (\wpt)^2$ with $C = 1/8$, compared against
        the \emph{sum}, never $\epsilon_E$ alone.} A resonantly driven Langmuir wave
        equipartitions, so the field holds half the wave energy on cycle average and
        swings between 3\% and 85\% of it instantaneously at $2\omega_p$.
  \item \textbf{The electron thermal reference is the 1V value
        $\epsilon_{\rm th} = \tfrac12 n_e \vthsq = 5\times10^{-4}$}, not
        $T00\_1 - 1 = 1.5\times10^{-3}$, which is the 3V rest-frame value.
  \item \code{Rho} in the field output is \emph{mass} density ($1 + 1836 = 1837$ in
        code units); the charge density is the separate \code{Charge} field.
  \item \code{CustomStat} receives \code{time} one step behind the \code{time} column
        of the same row (\code{engine.hpp} passes \code{time - dt} to the writer). This
        affects \code{e\_ext} bookkeeping only; the pusher's centring is correct.
\end{itemize}

\subsection{Reading the field dumps}

Field and particle output is ADIOS2 BP5 and is read with \code{nt2py}:

\begin{lstlisting}
pip install nt2py               # brings adios2 + xarray + dask
\end{lstlisting}
\begin{lstlisting}
import nt2
data = nt2.Data("<rundir>/<name>/<name>")     # note: name repeated
\end{lstlisting}

\code{pgens/dpdm/run34/extract\_run34.py} is the worked pipeline: it reduces a full
survey (${\sim}2.7$\,GB of \code{.bp}) to ${\sim}65$\,MB of \code{.npz}
--- \code{<tag>\_stats.npz}, \code{<tag>\_fields.npz}, \code{<tag>\_temps.npz},
plus \code{summary.npz} and \code{meta.json} --- so that plotting needs only
\code{numpy}. Set \code{run34\_dir} at the top of the file and run it once. This is the
intended route for analysing on a laptop: reduce on the cluster, download the
\code{.npz}.

\code{run34/PLOTS.md} states which four plots matter and what each one claims: the
energy budget (paper Fig.~2), the $E(k,t)$ and $n(k,t)$ spectra that show
\emph{why} it saturates, the $N_p$ convergence gate, and the fast/slow regime
transition. \code{run34\_plots.ipynb} produces them from the \code{.npz}.

\code{pgens/dpdm/paper/extract\_paper.py} is the same pipeline for the production
runs, which \code{run34} cannot substitute for: \code{run34} is a random-start survey
at a uniform $\Npc = 2000$ out to $\wpt = 3000$, whereas these are quiet-start, at the
paper's own $\Npc$, and include the swept-frequency family. It covers \code{fast},
\code{slow}, \code{lz1p0}, \code{heat} and the \code{undriven\_q}/\code{undriven\_r}
pair, adding $(x,u_x)$ phase-space histograms to the products, and reduces
${\sim}1.9$\,GB of \code{.bp} to ${\sim}50$\,MB. Stage the runs without their
checkpoint directories first --- \code{slow} alone carries 75\,GB of checkpoints
beside 483\,MB of science. \code{paper/PLOTS.md} states the six plots.

\section{Sizing a new run}
\label{sec:sizing}

Given a dark-photon amplitude $v_q^D/\vth$, in order:

\begin{enumerate}
  \item \textbf{Amplitude.} $A_0 = (v_q^D/\vth)\sqrt{10^{-3}}$.
  \item \textbf{Particle count.} Use the paper's shot-noise floor
        $\epsilon_{\rm noise} = L^2/(12 N_p N_x)$, in which \textbf{$N_p$ is the
        \emph{total} macroparticle count}, $N_p = \Npc N_x$ --- not the per-cell
        count. So evaluated, it agrees with measurement at a flat factor $3.1$ across
        three decades of $\Npc$, and is the right tool for this job. In code units the
        measured floor is $\epsilon_{\rm noise} = 2.06\times10^{-4}/\Npc$.
        \begin{itemize}[nosep]
          \item \emph{Resonant drive}: the signal grows as $t^2$, so it need only
                outgrow the noise before saturation.
          \item \emph{Swept drive}: the pre-resonance response is a \emph{bounded}
                forced oscillation, $A_0/(1-\omega^2/\omega_p^2) = 2.78A_0$ at
                $\omega = 0.8\omega_p$, so there is no time at which signal outgrows
                noise. Use the paper's criterion
                $\bar t_{\rm noise} = L/(A_0\sqrt{3N_pN_x/2}) \le 0.02$, giving
                $\Npc = 2.7\times10^3$ at $v_q^D/\vth = 1.0$ and $2.7\times10^5$
                at $0.1$.
        \end{itemize}
        \textbf{A quiet start does not reduce $\Npc$}; see \S\ref{sec:quietrelax}.
  \item \textbf{Memory, hence GPU count.} At $\Npc = 2.7\times10^5$ a run holds
        $5.3\times10^8$ particles and ${\approx}31$\,GB --- tight on one 40\,GB A100,
        comfortable at ${\approx}10.5$\,GB per GPU across three. Because the checkpoint
        domain count is locked (\S\ref{sec:restart}), settle this before submitting.
  \item \textbf{Walltime.} Take the per-step cost from Table~\ref{tab:perf}, then
        \textbf{double it}, or benchmark across the saturation transition. Per-step
        cost roughly doubles once the plasma becomes turbulent
        ($265 \to 516$\,ms on one A100 at $\Npc = 2\times10^5$); a benchmark run
        entirely inside the laminar phase measures the wrong number. Steps
        $= \text{runtime}/0.02$.
  \item \textbf{Queue.} \S\ref{sec:queues}. If the estimate exceeds the 48\,h cap,
        plan the restart chain from the start.
\end{enumerate}

\section{Rules and traps, consolidated}
\label{sec:rules}

Everything below is a standing rule. Each is stated once here; the evidence for it is
in Parts~\ref{part:impl} and~\ref{part:runs}.

\paragraph{Building.}
\begin{enumerate}[nosep]
  \item Configure and compile on a \textbf{login node}: \code{FetchContent} needs
        outbound network and compute nodes have none.
  \item \code{module load gcc/13.2.0} at build time \emph{and} in every job script.
  \item Pass \code{-D Kokkos\_ENABLE\_OPENMP=ON -D Kokkos\_ENABLE\_SERIAL=OFF}
        explicitly. \code{CMAKE\_IGNORE\_PATH}, which is required to make the ADIOS2
        fetch work, also blocks the repository's Kokkos settings, and a freshly
        fetched Kokkos defaults to Serial only --- a build that is externally
        indistinguishable and ${\sim}112\times$ slower. Neither SLURM's \code{-c}
        default nor \code{ibrun} binding causes this; do not go looking there.
  \item Always verify the cache after configuring
        (\S\ref{sec:quickstart}, step 2).
\end{enumerate}

\paragraph{Decks.}
\begin{enumerate}[nosep,resume]
  \item \code{ppc0} must be a TOML float.
  \item \code{loading = "quiet"} must be stated; the default is \code{"random"}.
        Do not use \code{"quiet\_lattice"} for new work --- it never decoheres
        (\S\ref{sec:quietrelax}).
  \item Include \code{[diagnostics]}, disable \code{[output.spectra]}, keep
        \code{[output.stats] interval\_time} $\ge 1.0$, keep
        \code{current\_filters = 0}.
\end{enumerate}

\paragraph{Running.}
\begin{enumerate}[nosep,resume]
  \item Pre-create \code{<rundir>/<simulation.name>/} before an MPI run.
  \item Fix the rank/GPU count before the first submission: checkpoints lock it.
  \item Budget a factor two on per-step cost across saturation.
  \item Jobs under 2\,h go to \code{development}; prefer fewer nodes.
  \item \code{blocking\_timers = true} is the first move on any performance surprise.
        It prints a per-kernel breakdown in ns per particle-step and settles in
        seconds questions that speculation gets wrong.
\end{enumerate}

\paragraph{Analysis.}
\begin{enumerate}[nosep,resume]
  \item Compare linear theory against $\epsilon_E + \Delta KE_e$, not $\epsilon_E$.
  \item Use the 1V thermal reference $5\times10^{-4}$.
  \item $N_p$ in the shot-noise formula is the total macroparticle count.
  \item \code{Rho} is mass density; charge density is \code{Charge}.
  \item $E_x(t{=}0) = 0$ proves nothing about the loading --- \textsc{Entity}
        initialises $E$ to zero identically and every loader co-locates the species.
        The discriminator is the total density $n_e + n_i$.
  \item A perfectly ordered lattice is \emph{not} the goal. The paper's initial
        condition is $\epsilon_{\rm noise}(0)=0$ \emph{followed by} relaxation to the
        Poisson floor within a few $\omega_p^{-1}$. Validate against both halves: a
        loader that stays ordered seeds daughter modes from round-off and delays
        saturation (\S\ref{sec:quietval}).
  \item \textsc{Entity}'s first output row is step~1, not $t=0$, so no dump shows the
        initial condition directly.
\end{enumerate}

\clearpage
\part{Implementation}
\label{part:impl}

\section{Extending \textsc{Entity}: this setup from scratch}
\label{sec:fromscratch}

This section assumes no prior knowledge of \textsc{Entity} and builds the DPDM
resonance setup from nothing. It is the orientation for the two sections that follow:
\S\ref{sec:pgen} describes what the production generator \emph{does},
\S\ref{sec:builds} how to compile it; this one explains what a problem generator
\emph{is}, how the pieces attach to the framework, and what the irreducible minimum
looks like.

\paragraph{The mental model.} \textsc{Entity}'s core is a fixed engine: Maxwell
solver, particle pusher, current deposit, filtering, I/O, MPI decomposition. None of
it is edited to add a new physics problem. Instead a \emph{problem generator}
(``pgen'') --- one C++ header --- is compiled into the engine, and it answers three
questions:

\begin{enumerate}[nosep]
  \item What particles exist at $t=0$, and where?
  \item What extra forces act on them, beyond the self-consistent fields?
  \item What extra numbers should be written to the diagnostics?
\end{enumerate}

Everything specific to this work is an answer to one of those three. The important
structural fact is that a binary contains \textbf{exactly one} problem generator,
chosen at configure time, not at run time --- so ``build the resonance case'' and
``build the Landau--Zener case'' are the same executable, while ``build the
reconnection case'' is a different one.

\subsection{Step 1: create the file}

\begin{lstlisting}
pgens/dpdm/pgen.hpp
\end{lstlisting}

That path is the entire registration procedure. CMake globs
\code{pgens/**/pgen.hpp} at configure time --- see \code{get\_available\_pgens}
in \code{cmake/config.cmake} --- and \code{-D pgen=dpdm} selects the directory whose
name matches. The selected header is then attached as an \textsc{interface} library
linked against the core targets, by \code{set\_problem\_generator} in the same
file. Three consequences worth stating plainly, because each removes a step one
would otherwise look for:

\begin{itemize}[nosep]
  \item there is no registry, enum or dispatch table in the core to add a name to;
  \item there is no \code{CMakeLists.txt} inside the pgen directory;
  \item the \emph{directory name} is the value of the \code{pgen} option --- there is
        no separate identifier declared inside the file.
\end{itemize}

Anything else in the directory (decks, plotting scripts, job scripts) is ignored by
the build and is there purely for the humans.

\subsection{Step 2: declare the generator and what it supports}

The header defines one class template in namespace \code{user}:

\begin{lstlisting}
template <SimEngine::type S, class M>
struct PGen {
  static constexpr auto D { M::Dim };
  static constexpr auto engines    = traits::pgen::compatible_with<SimEngine::SRPIC>{};
  static constexpr auto metrics    = traits::pgen::compatible_with<Metric::Minkowski>{};
  static constexpr auto dimensions = traits::pgen::compatible_with<Dim::_1D>{};
  ...
};
\end{lstlisting}

\code{S} is the physics engine and \code{M} the metric; the core instantiates the
template for whatever the deck asks for. The three \code{compatible\_with} traits are
declarations of intent, checked at compile time: they say this generator is written
for special-relativistic PIC on a flat Minkowski grid in one dimension. A deck asking
for anything else is rejected rather than silently producing wrong physics. Declaring
more than one value is allowed --- the production generator declares
\code{Dim::\_1D, Dim::\_2D, Dim::\_3D} --- but every declared combination has to
actually work.

\subsection{Step 3: read the deck}

The constructor receives the parsed input file and the global domain:

\begin{lstlisting}
PGen(const SimulationParams& p, const Metadomain<S, M>& global_domain)
  : params { p }
  , A0    { p.template get<real_t>("setup.A0", ZERO) }
  , omega { p.template get<real_t>("setup.omega", ONE) }
  , densities    { p.template get<prmvec_t>("setup.densities",    prmvec_t{TWO}) }
  , temperatures { p.template get<prmvec_t>("setup.temperatures", prmvec_t{ZERO,ZERO}) }
{ /* validation goes here */ }
\end{lstlisting}

This is where the \code{[setup]} block of \S\ref{sec:deck} comes from, and the
relationship is worth being explicit about: \textbf{\code{[setup]} is a free-form
namespace owned entirely by the problem generator.} The core never inspects it. There
is no schema, no list of legal keys, and nothing outside this one header knows that
\code{setup.drive} or \code{setup.loading} exist. Each option therefore has exactly
three touch points in \code{pgen.hpp} --- a member declaration, a \code{get} call with
its default, and the site that consumes it:

\begin{table}[h]
\centering\small
\begin{tabular}{@{}llll@{}}
\toprule
Key & declared & read (default) & consumed \\
\midrule
\code{setup.drive}       & \code{:301} & \code{:310} (\code{"resonance"})  & \code{drive\_phase()}, \code{:382} \\
\code{setup.loading}     & \code{:301} & \code{:311} (\code{"random"})     & \code{InitPrtls()}, \code{:433} \\
\code{setup.sweep\_phase}& \code{:303} & \code{:319} (\code{"integrated"}) & \code{drive\_phase()}, \code{:383} \\
\code{setup.A0}          & \code{:304} & \code{:321} (\code{0})            & \code{ext\_efield()}, \code{:393} \\
\bottomrule
\end{tabular}
\caption{Line numbers in \code{pgens/dpdm/pgen.hpp}. Adding an option means editing
these three places and nothing else.}
\end{table}

\paragraph{A trap that follows directly from this design.} The second argument to
\code{get} is a fallback, so an absent key is legal --- which means a \emph{misspelled
key silently takes the default}. Writing \code{loadng = "quiet"} runs the random
loader, with no error and no warning, and the run looks entirely normal. Only the
\emph{value} of a key is ever validated, never its name. Validation is manual, in the
constructor body:

\begin{lstlisting}
raise::ErrorIf(loading != "random" and loading != "quiet" and
                 loading != "quiet_lattice",
               "setup.loading must be `random`, `quiet`, or `quiet_lattice`", HERE);
\end{lstlisting}

The same mechanism is used for cross-checks the deck cannot express on its own: that
exactly two species are defined, and that their charges are equal and opposite.

\subsection{Step 4: put particles in the box}

\begin{lstlisting}
void InitPrtls(Domain<S, M>& domain) {
  arch::InjectUniformMaxwellians<S, M>(params, domain, densities[0],
                                       { temperatures[0], temperatures[1] }, { 1, 2 });
}
\end{lstlisting}

One call. \code{arch::} is \textsc{Entity}'s library of ready-made ``archetypes'';
this one fills the domain with a co-located electron--ion pair at the requested
density and temperatures, drawing positions uniformly at random. For a great many
problems this line is the whole of particle initialisation, and for the minimal
resonance setup it is sufficient.

It is also the exact point at which this work had to leave the beaten path. The
injector exposes two customisation seams --- the velocity distribution
(\code{EnrgDistClass}) and a density profile (\code{SpatialDistClass}, which sets
\emph{how many} particles land in a cell) --- but the sub-cell \emph{position} is
hard-coded as a random draw inside the kernel, with no seam at all. A quiet start is
by definition a statement about positions, so it cannot be expressed through the
archetypes and has to be a hand-written kernel. That is the subject of
\S\ref{sec:quiet}, and it is why \code{setup.loading} exists at all.

\subsection{Step 5: add the dark photon}

The drive is an external electric field that acts on the Lorentz force only; it is
never a source term in Maxwell's equations. It takes two objects. First a field
functor --- this is the part that runs \emph{on the GPU}, once per particle:

\begin{lstlisting}
template <Dimension D>
struct ExtFields {
  ExtFields(real_t e1) : e1 { e1 } {}
  Inline auto ex1(const coord_t<D>&) const -> real_t { return e1; }   // ignores x
private:
  const real_t e1;
};
\end{lstlisting}

and then the hook that hands one to the engine each step:

\begin{lstlisting}
auto ExternalFields(simtime_t t, spidx_t, const Domain<S, M>&) const
  -> std::pair<bool, ExtFields<M::Dim>> {
  return { A0 != ZERO, ExtFields<M::Dim>{ A0 * math::cos(omega * t) } };
}
\end{lstlisting}

The split matters for understanding what is cheap and what is not.
\code{ExternalFields} is called \textbf{on the host}, once per species per timestep
(\code{src/kernels/pushers/sr\_policies.h}); the functor it returns is then evaluated
\textbf{on the device}, once per particle. Because this drive is uniform in space, all
of the time dependence lives in the host-side scalar and the per-particle work is a
single register read --- the cosine is evaluated a few times per step, not $10^8$
times. The \code{bool} is a cheap switch: return \code{false} and the engine compiles
the external-field branch out of the push entirely.

\paragraph{How the engine knows the hook is there.} It is not registered, inherited or
declared anywhere. \textsc{Entity} tests the generator against C++20 concepts in
\code{src/global/traits/pgen.h} --- \code{HasExternalFields} asks whether a method of
that name and shape exists --- and compiles the call in if it does, or a no-op if it
does not. So every hook is opt-in by simply writing it, and omitting one is not an
error. The flip side is that a method with a slightly wrong signature is not a
compile error either: it is silently \emph{not a hook}, and the feature just does not
happen. The run log prints which hooks were detected (\code{engines/reporter.h}),
which is the quickest way to confirm one took effect.

\subsection{Step 6: report the drive (optional)}

\begin{lstlisting}
auto CustomStat(const std::string& name, timestep_t, simtime_t t,
                const Domain<S, M>&) const -> real_t {
  if (name == "e_ext") { return ext_efield(t); }
  return ZERO;
}
\end{lstlisting}

Named entries requested under \code{[diagnostics]} are routed here and land as extra
columns in the statistics CSV, alongside the field and particle energies. This is how
the drive's amplitude and phase become directly comparable with the plasma response in
the same file rather than having to be reconstructed afterwards. It is optional, and
nothing else depends on it.

\subsection{The minimal generator, in full}

Steps 2--5 assembled. This is a complete, working resonance setup:

\begin{lstlisting}
namespace user {
  using namespace ntt;
  using prmvec_t = std::vector<real_t>;

  template <Dimension D>
  struct ExtFields {
    ExtFields(real_t e1) : e1 { e1 } {}
    Inline auto ex1(const coord_t<D>&) const -> real_t { return e1; }
  private:
    const real_t e1;
  };

  template <SimEngine::type S, class M>
  struct PGen {
    static constexpr auto D { M::Dim };
    static constexpr auto engines    = traits::pgen::compatible_with<SimEngine::SRPIC>{};
    static constexpr auto metrics    = traits::pgen::compatible_with<Metric::Minkowski>{};
    static constexpr auto dimensions = traits::pgen::compatible_with<Dim::_1D>{};

    const SimulationParams& params;
    const real_t   A0, omega;
    const prmvec_t densities, temperatures;

    PGen(const SimulationParams& p, const Metadomain<S, M>&)
      : params { p }
      , A0    { p.template get<real_t>("setup.A0", ZERO) }
      , omega { p.template get<real_t>("setup.omega", ONE) }
      , densities    { p.template get<prmvec_t>("setup.densities",    prmvec_t{TWO}) }
      , temperatures { p.template get<prmvec_t>("setup.temperatures", prmvec_t{ZERO,ZERO}) } {}

    void InitPrtls(Domain<S, M>& domain) {
      arch::InjectUniformMaxwellians<S, M>(params, domain, densities[0],
                                           { temperatures[0], temperatures[1] }, { 1, 2 });
    }

    auto ExternalFields(simtime_t t, spidx_t, const Domain<S, M>&) const
      -> std::pair<bool, ExtFields<M::Dim>> {
      return { A0 != ZERO, ExtFields<M::Dim>{ A0 * math::cos(omega * t) } };
    }
  };
}
\end{lstlisting}

Roughly forty lines, of which two are physics. Everything in the production file
beyond this is either a loader (\S\ref{sec:quiet}) or a second drive family
(\S\ref{sec:pgen}).

\subsection{Step 7: a deck, a build, a run}

The deck needs the standard \textsc{Entity} sections --- \code{[simulation]},
\code{[grid]}, \code{[scales]}, \code{[algorithms]}, \code{[particles]},
\code{[diagnostics]}, \code{[output]}, covered in \S\ref{sec:deck} --- plus the
\code{[setup]} block this generator just defined:

\begin{lstlisting}
[setup]
  densities    = [2.0]
  temperatures = [1.0e-3, 1.0e-3]
  A0           = 9.48683e-4
  omega        = 1.0
\end{lstlisting}

Then configure, compile and run. \S\ref{sec:builds} is the full recipe with the three
traps that make the configure line longer than it looks; the short form on a login
node is:

\begin{lstlisting}
module load gcc/13.2.0
cmake -B build-dpdm -D pgen=dpdm -D precision=double \
      -D deposit=esirkepov -D shape_order=5 \
      -D Kokkos_ENABLE_OPENMP=ON -D Kokkos_ENABLE_SERIAL=OFF \
      -D Kokkos_ARCH_ZEN3=ON -D CMAKE_IGNORE_PATH=$PWD/cmake
cmake --build build-dpdm -j
grep -E "^Kokkos_ENABLE_(OPENMP|SERIAL|CUDA):BOOL" build-dpdm/CMakeCache.txt
./build-dpdm/src/entity.xc -input resonance.toml
\end{lstlisting}

\code{-D pgen=dpdm} is not optional: the default is \code{"."}, and with it CMake
configures the libraries but never adds the executable target, so the build
``succeeds'' and produces no \code{entity.xc}. The \code{grep} is the verification
step from \S\ref{sec:builds} --- a Serial build is indistinguishable from the outside
and runs ${\sim}112\times$ slower.

\subsection{Summary of the interface}

\begin{table}[h]
\centering\small
\begin{tabular}{@{}lll@{}}
\toprule
Hook & Detected by & Role here \\
\midrule
\emph{(constructor)}  & required           & reads \code{[setup]}, validates it \\
\code{engines}, \code{metrics}, \code{dimensions} & required & declares what the deck may ask for \\
\code{InitPrtls}      & \code{HasInitPrtls}      & the $t=0$ particle state \\
\code{ExternalFields} & \code{HasExternalFields} & the dark-photon drive \\
\code{CustomStat}     & \code{HasCustomStatOutput}     & writes \code{e\_ext} to the CSV \\
\bottomrule
\end{tabular}
\caption{Everything below the first two rows is opt-in: a generator that defines none
of them compiles and runs, it simply has no particles and no extra physics. Concepts
are in \code{src/global/traits/pgen.h}.}
\end{table}

What the production generator adds on top of the minimum, and where each is treated:
the deterministic loader and its two variants (\S\ref{sec:quiet}); the swept
Landau--Zener phase (\S\ref{sec:pgen}) and the two conventions for it
(\S\ref{sec:sweep}); the 1V option, which zeroes the two transverse velocity
components to match \textsc{sharp}; and the \code{e\_ext} diagnostic.

\paragraph{On naming options rather than numbering them.} Two of the four
\code{[setup]} switches are, strictly, redundant. \code{drive = "resonance"} is
identical to \code{"landau\_zener"} with $\dot\omega=0$, since
$\phi = \omega_0 t + \tfrac12\dot\omega t^2$ collapses to $\omega_0 t$; and
\code{sweep\_phase = "naive"} is \emph{exactly} the integrated form at twice the
sweep rate, since $(\omega_0 + \dot\omega t)t = \omega_0 t +
\tfrac12(2\dot\omega)t^2$. The whole drive could therefore be three numbers ---
$A_0$, $\omega_0$, $\dot\omega$ --- with no strings and no validation branches at
all. The names are kept because a deck reading \code{domega\_dt = 4.0e-5} records
what was run but not \emph{why}, and for a reproduction the intent is part of the
result; the strings also propagate into output paths and run names. It is a
legibility choice, not a necessity, and a generator written for a single case should
probably not make it.


\section{Problem generator}
\label{sec:pgen}

All problem-specific code lives in \code{pgens/dpdm/pgen.hpp}. It provides three
things.

\paragraph{External field.} An \code{ExternalFields} hook returning a uniform
$e_{x1}$, applied to \emph{both} species; the pusher scales by $q/m$ per species, so
the ions automatically receive $1/1836$ of the electron kick. Because the drive is
uniform in space, the cosine is evaluated once per species per step on the host and
only the scalar is handed to the kernel --- not once per particle.

\paragraph{Two drive families.} Selected by \code{setup.drive}:
\begin{align}
  \text{resonance:}\quad & \phi(t) = \omega t, \qquad \omega = \omega_p
  \ \text{fixed},\\
  \text{Landau--Zener:}\quad & \phi(t) = \omega_0 t + \tfrac12 \dot\omega t^2 .
\end{align}
The swept phase is integrated \emph{analytically}. Sampling $\omega(t)$ per step and
multiplying by $t$ accumulates a spurious chirp: the naive form deviates by up to
$0.86\,A_0$ over the sweep. The analytic form was verified against the closed
expression to $2.8\times10^{-14}$ (double precision, i.e.\ exact).

\paragraph{A custom diagnostic.} \code{CustomStat} exposes \code{e\_ext}, the
instantaneous drive, so amplitude and phase can be checked against the plasma response
in the same CSV. Note that \code{CustomStat} receives \code{time} one timestep behind
the \code{time} column in the same row (\code{engine.hpp} passes \code{time - dt} to
the writer); this is a diagnostic bookkeeping offset only --- the pusher's centring is
correct, as the drive-normalisation measurement in \S\ref{sec:drivenorm} confirms.

\section{Deterministic (``quiet'') particle loading}
\label{sec:quiet}

The paper initialises ``all particles equally spaced, with ions and electrons
colocated throughout the computational domain,'' which ensures
$\epsilon_{\rm noise}=0$ at $t=0$. \textsc{Entity}'s stock injector places every
particle at a uniformly random position in the domain --- not even one per cell ---
so cell occupancy is Poisson and $\epsilon_{\rm noise}(0)\neq 0$. A loader was
therefore written for this work.

\subsection{What the initial condition actually requires}
\label{sec:quietreq}

Two things, and the second is easy to miss.

\begin{enumerate}[nosep]
  \item \textbf{$\epsilon_{\rm noise}=0$ at $t=0$.} Stated explicitly in the paper.
  \item \textbf{Prompt relaxation to the Poisson floor.} The paper's Fig.~A00 shows
        undriven runs starting at zero and rising to the analytic level
        $L^2/(12N_pN_x)$ within a few $\omega_p^{-1}$, then sitting flat for
        $2\times10^4$. The text says the same from the driven side: ``it takes
        approximately $t\omega_p \sim 3$--$4$ for the driven energy to be converted
        into de-coherent (high-$k$) noise.''
\end{enumerate}

The second requirement is not decorative. The paper's entire $\Npc$ criterion
($\bar t_{\rm noise}\sim 2.4$ resonant, $\le 0.02$ swept) presumes that shot noise
is present and sets the seed from which daughter modes grow. A loader that
satisfies only the first produces a plasma with \emph{no} seed at all.

\subsection{Three noise channels}

\begin{enumerate}[nosep]
  \item \textbf{Charge density.} $\rho(x,0)\neq 0$ seeds $E(x,0)\neq0$. Killed by
        co-locating each electron with an ion. \textsc{Entity}'s random start
        already does this, so this channel was never the problem.
  \item \textbf{Total density.} $n_{\rm tot}=n_e+n_i$ fluctuates at
        $\sqrt{2/\Npc}$ even with perfect co-location. Since
        $\omega_p\propto\sqrt{n}$ this is a spread in the local resonant frequency
        --- the channel that physically matters.
  \item \textbf{Velocity space.} Randomly sampled velocities produce current
        fluctuations that regenerate $E$ within a plasma period.
\end{enumerate}

Channel 2 must be \emph{zero at $t=0$} and must \emph{reappear at the Poisson level
within a few $\omega_p^{-1}$}.

\subsection{Construction}
\label{sec:quietconstr}

Positions form an exact lattice --- cell $i$ holds $\Npc$ particles at sub-cell
offsets $(s+\tfrac12)/\Npc$ --- and ions reuse them exactly, so both channel~1 and
channel~2 vanish to round-off at $t=0$. Velocities are drawn pseudo-randomly,
keyed on a \code{splitmix64} hash of the \emph{global} cell index and slot, which
keeps the loading independent of how MPI decomposes the domain.

\begin{lstlisting}
i   = p / ppc                         // cell
s   = p % ppc                         // slot within the cell
phi = (s + 1/2)/ppc                   // exact lattice: eps_noise(0) = 0
x   = xi_min + i + phi
key = (global cell)*1000003 + s       // decomposition independent
u_e = sig1 * probit(hash_u01(4*key))
u_i = sig2 * probit(hash_u01(4*key+1))
\end{lstlisting}

\paragraph{Why the velocity assignment must be irregular.} Two constructions that
look reasonable both fail, for the same reason.

\emph{Velocity keyed on the class.} Give the same velocity to the same slot in
every cell. Each velocity class is then a comb of exactly one particle per cell,
spacing $\Delta x$, all moving at one speed. A uniform comb translates rigidly and
deposits exactly uniform density at any offset, because the shape function is a
partition of unity at that spacing. The sum of such combs is uniform forever: the
lattice \textbf{never decoheres}.

\emph{A rigid per-cell rotation.} Shifting the assignment by $iR$ in cell $i$ does
not help --- each class is still a uniform comb, now of spacing
$(1-R/\Npc)\Delta x$. Still rigid, still no clumping.

The assignment has to be irregular \emph{from cell to cell}, which is what the hash
provides. A worked example with $N_x=4$, $\Npc=4$ is the quickest way to see it: with
class-keyed velocities, class $A$ sits at $x = 0.125, 1.125, 2.125, 3.125$ --- spacing
exactly $1$, a perfect comb. With hashed velocities it sits at
$x = 0.375, 1.125, 2.875, 3.625$ --- spacings $0.75, 1.75, 0.75$, irregular, and it
clumps as soon as it streams. Both have identical, perfectly flat density at $t=0$.

\paragraph{Velocity normalisation, and why it is now skipped.} An ordered stratified
set $\{\mathrm{probit}((j+\tfrac12)/\Npc)\}$ has mean zero but variance below unity
($-1.97\%$ at $\Npc=64$, $-0.64\%$ at $200$), so the legacy loader measures the
realised variance and divides it out. The sampling loader must \emph{not} do this:
its realised variance carries a physical $1/\sqrt{N}$ fluctuation, the velocity-space
counterpart of the shot noise being restored, and normalising it away would re-quiet
the channel the loader exists to open.

\subsection{Validation: the loader gate}
\label{sec:quietval}

Three undriven runs ($A_0=0$, $\Npc = 2700$, identical decks) differing only in
\code{setup.loading}. Analytic floor $L^2/(12N_pN_x) = 4.938\times10^{-8}$; the table
gives $\epsilon_E$ as a fraction of it.

\begin{table}[h]
\centering\small
\begin{tabular}{@{}lrrrrr@{}}
\toprule
\code{loading} & $\wpt=1$ & $10$ & $50$ & $200$ & $500$ \\
\midrule
\code{quiet} (sampling)        & 1.08 & 1.33 & 1.27 & 1.45 & 1.34 \\
\code{quiet\_lattice} (legacy) & $5.1\times10^{-14}$ & $2.1\times10^{-12}$ &
                                 $1.7\times10^{-11}$ & $2.1\times10^{-10}$ &
                                 $3.4\times10^{-10}$ \\
\code{random} (stock)          & 1.18 & 1.54 & 1.53 & 1.44 & 1.58 \\
\bottomrule
\end{tabular}
\caption{Loader gate, job 3447505. \code{quiet} is at the Poisson floor by
$\wpt=1$ and flat --- Fig.~A00. The legacy loader is 10--15 orders below it and
still climbing at $\wpt=500$. \code{random} is the control: it lands on the same
floor, confirming the formula and the $N_p$ convention independently of any
quiet-start logic. Script: \code{run\_gate\_loader.sh}.}
\end{table}

Two traps. \textsc{Entity}'s first output row is \textbf{step 1}, not $t=0$, so no
dump shows the initial condition directly; $t=0$ uniformity is exact by construction,
and \code{quiet\_lattice} measuring $\delta n/n = 0.000$ with
$\max|\rho_q| = 6\times10^{-17}$ at step~1 confirms the deposition is exact for that
position set. And $\epsilon_E(0)\approx 0$ for \emph{all three} loadings, including
\code{random}, because \textsc{Entity} initialises $E$ to zero identically --- the
discriminator is the total density, never the field.

\subsection{Consequences of the legacy loader}
\label{sec:quietrelax}

The legacy loader was used for the first round of production runs, and their
saturation times and peak energies are not reproductions. Measured $k\neq0$ field
power relative to the $k=0$ driven mode at $\wpt\approx 40$:

\begin{center}\small
\begin{tabular}{@{}lrr@{}}
\toprule
Run & $P_{k\neq0}/P_{k=0}$ & $\wpt$ at which the ratio reaches $1\%$ \\
\midrule
\code{fast}, legacy    & $2.7\times10^{-31}$ & 962 \\
\code{fast}, sampling  & $1.0\times10^{-2}$  & 0 \\
\code{lz1p0}, legacy   & $1.7\times10^{-29}$ & 1120 \\
\code{lz1p0}, sampling & $1.1\times10^{-3}$  & 0 \\
\bottomrule
\end{tabular}
\end{center}

$2.7\times10^{-31}$ is double-precision epsilon \emph{squared}: the daughter modes
were growing out of round-off. At $\gamma = 0.061\,\omega_p$, climbing from $10^{-31}$
to order unity takes ${\sim}1100\,\omega_p^{-1}$ of extra runway, during which the
drive keeps pumping.

\textbf{The effect is large where saturation time sets the amplitude, and small
otherwise.} In the resonant cases the $k=0$ wave grows as $t^2$ until daughters halt
it, so a delay translates directly into a larger peak. In the swept case the energy
available is set by how long $\omega$ lingers near $\omega_p$, not by the onset of
turbulence, so the same delay changes little:

\begin{center}\small
\begin{tabular}{@{}lrrr@{}}
\toprule
Run & peak $\epsilon_E/\epsilon_{\rm th}$ & $\wpt$ at peak & $kT_e$(end)$/kT_e(0)$ \\
\midrule
\code{fast}  legacy $\rightarrow$ sampling & $43.55 \rightarrow 9.14$ & $553 \rightarrow 238$ & $70.7 \rightarrow 75.5$ \\
\code{slow}  legacy $\rightarrow$ sampling & $2.13 \rightarrow 0.53$  & $3694 \rightarrow 1598$ & $3.43 \rightarrow 2.53$ \\
\code{lz1p0} legacy $\rightarrow$ sampling & $190.4 \rightarrow 237.6$ & $3689 \rightarrow 4269$ & $1003 \rightarrow 962$ \\
\bottomrule
\end{tabular}
\end{center}

\paragraph{A quiet start still does not permit a reduction in $\Npc$.} With the
sampling loader $\Npc$ sets the seed amplitude directly, which is what the paper's
sizing assumes; under the legacy loader round-off set it and $\Npc$ was doing nothing
for the physics. Either way, particle counts must be sized from the shot-noise
criterion (\S\ref{sec:sizing}).

\section{Building on Lonestar6}
\label{sec:builds}

This section is the full recipe: the four build directories used in this work can be
reconstructed from scratch by following it in order. The directories themselves are
not version-controlled --- they total ${\sim}1.2$\,GB of object files, fetched
dependency sources and architecture-specific binaries, and are valid only on this
machine.

\subsection{Environment and prerequisites}

All four builds are configured and compiled \textbf{on a login node}. \textsc{Entity}
fetches Kokkos and ADIOS2 with CMake's \code{FetchContent} at configure time, and LS6
compute nodes have no outbound network: the same commands on a compute node fail in
the clone step.

\begin{lstlisting}
module load gcc/13.2.0    # switches the hierarchy; brings impi/21.12 with it
module load cuda/12.8     # CUDA builds only -- LS6 defaults to cuda/11.4
\end{lstlisting}

\code{gcc/13.2.0} is not optional. The binaries link against GCC~13's
\code{libstdc++}; the default LS6 environment provides gcc/9.4.0, under which the
executable dies at load time with \code{GLIBCXX\_3.4.31 not found} \emph{before} it
reads its input file. The module must therefore also be loaded in every job script,
not only at build time. Loading it switches the module hierarchy to
\code{/opt/apps/gcc13\_2/modulefiles} and pulls in \code{impi/21.12} (Intel oneAPI MPI
2021.12, \code{mpigxx}), which is the MPI the two parallel builds pick up. CUDA must be
requested explicitly: the site default \code{cuda/11.4} is too old for the fetched
Kokkos, and the builds here used \code{cuda/12.8} (\code{nvcc} at
\code{/opt/apps/cuda/12.8/bin/nvcc}).

\subsection{Configure}

Four configurations are used, differing only in backend. All are configured from the
repository root:

\begin{lstlisting}
IGN=-DCMAKE_IGNORE_PATH=$PWD/cmake
COMMON="-D pgen=dpdm -D deposit=esirkepov -D shape_order=5 -D precision=double"
KOK="-D Kokkos_ENABLE_OPENMP=ON -D Kokkos_ENABLE_SERIAL=OFF"

# 1. CPU, OpenMP           -- validation, fast resonance case
cmake -B build-dpdm          $COMMON $KOK -D Kokkos_ARCH_ZEN3=ON $IGN
# 2. CPU + MPI             -- multi-node fallback
cmake -B build-dpdm-mpi      $COMMON $KOK -D Kokkos_ARCH_ZEN3=ON -D mpi=ON $IGN
# 3. CUDA, single GPU      -- the production path
cmake -B build-dpdm-cuda     $COMMON -D Kokkos_ENABLE_CUDA=ON \
      -D Kokkos_ARCH_AMPERE80=ON $IGN
# 4. CUDA + MPI, multi-GPU -- LZ 0.1, long-term stability
cmake -B build-dpdm-cuda-mpi $COMMON $KOK -D mpi=ON -D gpu_aware_mpi=OFF \
      -D Kokkos_ENABLE_CUDA=ON -D Kokkos_ARCH_AMPERE80=ON $IGN
\end{lstlisting}

\code{Kokkos\_ARCH\_ZEN3} targets the AMD EPYC 7763 CPU nodes and
\code{Kokkos\_ARCH\_AMPERE80} the A100s; configuration~3 omits \code{\$KOK} because a
CUDA backend does not need a threaded host backend for this workload. Each \code{cmake
-B} re-fetches and re-builds its own private copy of Kokkos and ADIOS2 under
\code{<build>/\_deps}, which is why four directories cost ${\sim}1.2$\,GB rather than
four small binaries.

\paragraph{Historical note.} Configuration~1 was in fact first configured with
\code{-D adios2\_DIR=\$PWD/build/\_deps/adios2-build}, reusing the ADIOS2 that an
earlier unrelated \code{build/} directory had already fetched. That works, but it is
not reproducible from a clean checkout --- there is no \code{build/} to point at --- so
the form given above, using \code{\$IGN} like the other three, is the one to follow.
The two are equivalent: both end up at ADIOS2 2.11.0.

\subsection{Compile}

\begin{lstlisting}
cmake --build build-dpdm          -j 8
cmake --build build-dpdm-mpi      -j 8
cmake --build build-dpdm-cuda     -j 8
cmake --build build-dpdm-cuda-mpi -j 8
\end{lstlisting}

Each produces its executable at \code{<build>/src/entity.xc}; that path is what the job
scripts in \code{pgens/dpdm/} invoke. The CUDA configurations take substantially longer
than the CPU ones --- every kernel is compiled twice, to PTX and to Ampere SASS.
Keep \code{-j} modest: LS6 login nodes are shared, and a large \code{-j} on an
\code{nvcc} build is the usual way to get a job killed by the login-node watchdog.

\paragraph{Verify before trusting a build.} Configure-time options fail silently here
(Trap~2 below), so check the cache rather than assuming:
\begin{lstlisting}
grep -E "^Kokkos_ENABLE_(OPENMP|SERIAL|CUDA):BOOL" build-*/CMakeCache.txt
\end{lstlisting}

\paragraph{Trap 1: the ADIOS2 shim shadows its own package.}
\code{find\_package(adios2 QUIET)} matches the repository's own
\code{cmake/adios2Config.cmake}, which only sets cache variables and exports neither a
version nor any targets. CMake nevertheless sets \code{adios2\_FOUND} and the fetch
branch never runs. The symptom is \code{ADIOS2: v} with a path pointing into
\code{entity/cmake}, followed by \code{Target "entity.xc" links to adios2::cxx\_mpi but
the target was not found}. Excluding that directory via \code{CMAKE\_IGNORE\_PATH}
restores the intended fetch. This is an upstream bug: the shim shadows the package it
exists to configure.

\paragraph{Trap 2: \code{CMAKE\_IGNORE\_PATH} also disables Kokkos OpenMP.}
The same exclusion blocks the repository's Kokkos settings file, and a freshly fetched
Kokkos defaults to \textbf{Serial only}. A Serial build is externally
indistinguishable from a threaded one and runs ${\sim}112\times$ slower
(\code{NLWP = 2}, 99.5\% CPU). Neither SLURM's \code{-c} default nor \code{ibrun}
binding produces this symptom, so there is no point looking there; the cache check at
the end of the previous subsection is mandatory rather than advisory.

\paragraph{Trap 3: ADIOS2 target rename.} ADIOS2~2.11 renamed the C++ bindings target
from \code{cxx11} to \code{cxx} (and \code{cxx11\_mpi} to \code{cxx\_mpi}).
\textsc{Entity} hardcodes the 2.11 names, so TACC's \code{adios2/2.10.2} module cannot
link even though it does provide \code{ADIOS2\_HAVE\_MPI}. Fetching 2.11.0
(\textsc{Entity}'s pinned tag) is the clean route, and requires a login node --- compute
nodes have no outbound network.

\paragraph{GPU-aware MPI.} Disabled (\code{gpu\_aware\_mpi=OFF}). It would require
\code{mvapich2-gdr-cuda12}, a different MPI stack from the \code{impi} in the gcc13
hierarchy. The halo here is a few ghost cells of three field components, so host
staging is free at this problem size.

\section{Measured performance}
\label{sec:perf}
\label{sec:running}

This section records the throughput measurements that the sizing procedure of
\S\ref{sec:sizing} draws on.

\subsection{Throughput}

\begin{table}[h]
\centering\small
\begin{tabular}{@{}lrl@{}}
\toprule
Configuration & ms/step at $\Npc = 2\times10^5$ & Note \\
\midrule
1 CPU node (128 cores)     & ${\sim}950$ & steady state; rises $+48\%$ from 640 \\
8 CPU nodes, MPI           & 135         & $4.94\times$ speedup, 69\% efficiency \\
$1\times$ A100             & $265 \rightarrow 516$ & $3.6\times$ a full CPU node \\
$3\times$ A100 (one node)  & 128         & $2.07\times$ over 1 GPU, 69\% efficiency \\
\bottomrule
\end{tabular}
\caption{Measured on $4\times10^8$ particles, $N_x = 1000$.}
\label{tab:perf}
\end{table}

A single A100 outperforms an entire 128-core node by $3.6\times$. The feared collapse
from deposition atomics on a 1000-cell grid did not occur: going from
$\Npc = 2\times10^4$ to $2\times10^5$ --- a tenfold increase in particles per cell, and
therefore in atomic contention per cell --- cost only a factor $1.42$
($0.444 \rightarrow 0.630$\,ns per particle-step).

\subsection{Per-step cost doubles after saturation}
\label{sec:slowdown}

This governs all walltime planning.

On a single A100 at $\Npc = 2\times10^5$ the cost is $265$\,ms/step while the plasma is
laminar and $516$\,ms/step once it saturates and becomes turbulent --- particles
disorder, and the memory locality the early rate depended on is lost. A 2000-step
benchmark ($\wpt = 40$) shows a dead-flat $265$\,ms, which is \emph{not} evidence of
stability: the run saturates at $\wpt \approx 2000$, so such a benchmark never leaves
the laminar phase. Taken at face value it overran a 48\,h wall at 88\% completion on
the slow resonance case, which then required a checkpoint restart.

\textbf{Rule: benchmark across the physics transition, or budget a factor two.}

\clearpage
\part{Runs, settings, and comparison with the paper}
\label{part:runs}

\section{Parameters taken from the paper source}
\label{sec:paperparams}

All values below are quoted from \code{main.tex} in the arXiv source package, not
inferred from the text or figures.

\begin{table}[h]
\centering\small
\begin{tabular}{@{}lll@{}}
\toprule
Quantity & Paper value & This work \\
\midrule
Grid                  & $N_x = 1000$, $L = 40\,c/\omega_p$ & identical \\
Cell size             & $\Delta x = 0.04\,c/\omega_p$      & identical \\
Debye length          & $\lambda_D \approx 0.033\,c/\omega_p$ & identical \\
Temperatures          & $k_BT_e = k_BT_i = 10^{-3}m_ec^2$  & identical \\
Mass ratio            & $m_i/m_e = 1836$                   & identical \\
Thermal speed         & $\vthsq \equiv \langle(v_x-\langle v_x\rangle)^2\rangle = k_BT/m$ & identical (1V) \\
Thermal energy density& $5\times10^{-4}\,n_em_ec^2$        & $4.9963\times10^{-4}$ measured \\
Deposition            & 5th-order splines                  & \code{esirkepov}, \code{shape\_order=5} \\
Pusher                & Vay (2008)                         & \code{Vay} \\
Timestep              & $\Delta t \approx 0.02/\omega_p$   & CFL $=0.5 \Rightarrow 0.02$ \\
Initial loading       & equally spaced, co-located         & \S\ref{sec:quiet} \\
Shot noise            & $\epsilon_{\rm noise} = L^2/(12 N_p N_x)$ & agrees to $3.1\times$ \\
Resonance $\Npc$      & $2\times10^2$ and $2\times10^5$    & identical \\
LZ sweep              & $\omega(t) = 0.8\omega_p + 0.1\omega_p t/5000$ & identical \\
LZ amplitudes         & $v_q^D/\vth = 1.0$ and $0.1$       & identical \\
LZ $\Npc$ criterion   & $\bar t_{\rm noise} \le 0.02$      & $2.7\times10^3$, $2.7\times10^5$ \\
\bottomrule
\end{tabular}
\caption{Paper parameters and the corresponding settings used here.}
\end{table}

Two points deserve emphasis.

\paragraph{The 1V definition.} The paper defines the thermal speed from the
\emph{one-dimensional} $x$-component variance, $\vthsq = \langle (v_x - \langle v_x
\rangle)^2\rangle = k_BT/m$. \textsc{Entity}'s Maxwellian is isotropic 3V. Because
there are no transverse forces in a 1D electrostatic run, the perpendicular degrees of
freedom merely carry constant energy; including them would dilute $T_e$ by a factor
three. The quiet loader therefore sets $u_{x2}=u_{x3}=0$ (\code{one\_v = true}),
matching \textsc{SHARP} exactly and removing the 3V-versus-1V correction from every
downstream comparison. The paper's quoted thermal energy density $5\times10^{-4}$
matches the 1V value and is used as the reference line throughout.

\paragraph{The Landau--Zener crossing time.} With $\omega_0 = 0.8$ and
$\dot\omega = 2\times10^{-5}$, resonance ($\omega = \omega_p$) is crossed at
$\wpt = 10^4$. The paper's statement that ``nonlinearity kills the resonant transfer
around $\wpt = 5000$'' refers to $\omega = 0.9\,\omega_p$, i.e.\ \emph{before} the
crossing --- off-resonant excitation, not the crossing itself.

\section{Preparatory and calibration runs}

Before the production runs, a sequence of measurements pinned the unit conventions and
validated the drive. These are summarised here because several production choices
depend on them.

\begin{itemize}[nosep]
  \item \textbf{Units.} \code{skindepth0 = 1} with \code{densities = [2.0]} gives
        $n_e = n_0$ exactly, hence $\omega_{pe} = 1$. Confirmed by \code{T00\_1}
        rest-mass part $= 1.0000$.
  \item \textbf{Energy conservation.} Over $2\times10^4$ steps,
        $\Delta T^{00} = -\tfrac12\Delta\langle E^2\rangle$ with $r = -0.996$; no grid
        heating.
  \item \textbf{$\omega_p$ calibration.} A five-point detuning scan
        ($\omega = 0.96$--$1.04$) is symmetric about $\omega = 1$ and follows the
        linear-theory $\mathrm{sinc}^2$ envelope; the on-resonance case grows as $t^2$
        across two decades in energy while every detuned case remains bounded.
  \item \textbf{Drive normalisation.}\label{sec:drivenorm} Least squares over
        $10 < \wpt < 40$ gives $C = 0.12432$ against Eq.~A25's $1/8 = 0.125$
        (0.5\%). The ion response ratio is $5.70\times10^{-4}$ against
        $1/1836 = 5.45\times10^{-4}$, confirming the drive reaches both species with
        the correct $q/m$.
\end{itemize}

\paragraph{A factor-of-two trap.} It is the \emph{sum} $\epsilon_E + \Delta KE_e$ that
obeys the $t^2$ law, not the field energy alone. The resonantly driven Langmuir wave
equipartitions, so the field carries only half the wave energy on cycle average and
swings between 3\% and 85\% of it instantaneously at $2\omega_p$. The paper's
linear-theory line must be compared against the sum.

\section{The two resonance cases (paper Fig.~2)}
\label{sec:resonance}

Both run to $\wpt = 10^4$, 1V, with the sampling loader of \S\ref{sec:quietconstr}.
The legacy-loader values are shown alongside because they appeared in earlier drafts
and are \emph{not} reproductions (\S\ref{sec:quietrelax}).

\begin{table}[h]
\centering\small
\begin{tabular}{@{}llrrrr@{}}
\toprule
Case & loader & $\Npc$ & $A_0$ & peak $\epsilon_E/\epsilon_{\rm th}$ & $kT_e$(end)$/kT_e(0)$ \\
\midrule
fast, $v_q^D/\vth = 3\times10^{-2}$ & legacy   & $2\times10^2$ & $9.48683\times10^{-4}$ & 43.55 (at 553)  & 70.7 \\
                                    & sampling & $2\times10^2$ & $9.48683\times10^{-4}$ & \textbf{9.14} (at \textbf{238}) & \textbf{75.5} \\
slow, $v_q^D/\vth = 1\times10^{-3}$ & legacy   & $2\times10^5$ & $3.16228\times10^{-5}$ & 2.13 (at 3694) & 3.43 \\
                                    & sampling & $2\times10^5$ & $3.16228\times10^{-5}$ & \textbf{0.53} (at \textbf{1598}) & \textbf{2.53} \\
\bottomrule
\end{tabular}
\caption{Production resonance runs, legacy versus sampling loader.
$\epsilon_{\rm th} = 5\times10^{-4}$ is the 1V electron thermal energy density.
Runs \code{fast\_s} and \code{slow\_s}; the legacy rows are \code{fast}, \code{slow}.}
\end{table}

\paragraph{Both cases now agree with the paper.} For the fast-growth regime the paper
says the electron temperature grows to $O(10^2)$ of its initial value: we measure
$75.5$. For the slow-growth regime it says the transfer slows ``\emph{before} the
energy density stored in the Langmuir wave becomes comparable to the initial thermal
energy'' and that ``the total energy transfer is about the electron thermal energy'':
we measure a field peak of $0.53\,\epsilon_{\rm th}$, which \textbf{never crosses the
thermal line at any sample}, and a peak total transfer
$(\epsilon_E + \Delta KE_e)$ of $1.59\,\epsilon_{\rm th}$.

\paragraph{The regime distinction.} The paper separates slow and fast growth at
$v_q^D/\vth \sim (m_e/m_p)^{1/2}/2 \approx 0.0117$. The discriminator is the lasting
electron heating --- $75.5$ against $2.53$, a factor ${\sim}30$ --- not the peak field
energy. With the sampling loader the two cases also separate cleanly on whether the
field energy reaches the thermal line at all: the fast case exceeds it by $9\times$,
the slow case never reaches it.

\paragraph{Initial temperature.} $k_BT_e(0) = 4.9963\times10^{-4}$ against the 1V
target $5.0\times10^{-4}$. The $-0.075\%$ residual is \emph{not} loading error but the
relativistic definition $T^{00} = \sum w\,m\,\gamma$, whose next term is
$-\tfrac38\langle u^4\rangle = -0.075\%$. The ions, 1836 times heavier and far more
non-relativistic, land at $4.999998\times10^{-4}$, which confirms the interpretation.

\paragraph{Cost.} On the legacy loader the slow case took ${\sim}51$ GPU-hours on one
A100 (48\,h to 88\%, then a 14.5\,h restart). The sampling re-run used three A100s on
one node and took \textbf{43\,h} against an estimate of ${\sim}25$\,h from the measured
$2.07\times$ scaling: saturation arrives earlier with a physical seed, so more of the
run sits in the post-saturation regime where the per-step cost roughly doubles
(\S\ref{sec:slowdown}). Estimates built on the laminar rate are optimistic by about a
factor two even after that effect is known about.

\section{Landau--Zener runs}
\label{sec:lz}

The Landau--Zener family is the physically realistic case: in cosmology it is
$\omega_p(t)$ that drifts through the dark-photon mass. The paper instead sweeps
$m_{A'}(t)$ at a numerically tractable rate, quoting
$\omega_p/H \approx 7\times10^4$.

\subsection{The sweep convention}
\label{sec:sweep}

The paper's appendix gives the drive as $A_0\cos(\omega t)$ with
$\omega(t) = 0.8\omega_p + 0.1\omega_p t/5000$. That admits two readings, and they
are not equivalent:

\begin{center}\small
\begin{tabular}{@{}lllr@{}}
\toprule
\code{setup.sweep\_phase} & phase & instantaneous frequency & crossing \\
\midrule
\code{integrated} (default) & $\omega_0 t + \tfrac12\dot\omega t^2$ & $\omega_0+\dot\omega t$ & $\wpt = 10^4$ \\
\code{naive}                & $(\omega_0+\dot\omega t)\,t$          & $\omega_0+2\dot\omega t$ & $\wpt = 5000$ \\
\bottomrule
\end{tabular}
\end{center}

\code{integrated} is the correct realisation of ``a drive whose frequency evolves as
$\omega(t)$'' and is the default. \code{naive} is a well-defined chirp sweeping at
twice the nominal rate, provided so the paper's convention can be tested.

The paper's own evidence is contradictory. Its stated $\omega_p/H\approx7\times10^4$
matches \code{integrated}: with $\mathrm{d}\ln\omega_p^2/\mathrm{d}t = 3H$ and
$H = 2\dot\omega/(3\omega)$ at crossing, $\dot\omega = 2\times10^{-5}$ gives
$7.5\times10^4$, against $3.75\times10^4$ for the naive form. But its
\emph{figure} points the other way: the $x$-axis runs $0$ to $10^4$ and the
linear-theory curve completes its rise and goes flat by $\wpt\approx5300$. Linear
Landau--Zener theory transfers at the crossing, so the paper's crossing is near
$5000$ --- where \code{naive} crosses, and where \code{integrated} has nothing
happening. The figure is the stronger evidence, which makes the quoted
$\omega_p/H$ too large by a factor two.

\subsection{Results, and an unresolved discrepancy}

\begin{table}[h]
\centering\small
\begin{tabular}{@{}llrrr@{}}
\toprule
Run & loader / sweep & crossing & peak $\epsilon_E/\epsilon_{\rm th}$ & $kT_e$(end)$/kT_e(0)$ \\
\midrule
\code{lz1p0}   & legacy / integrated   & $10^4$ & 190.4 (at 3689) & 1003 \\
\code{lz1p0\_s}& sampling / integrated & $10^4$ & 237.6 (at 4269) & 962 \\
\code{lz1p0n}  & sampling / naive      & $5000$ & 246.7 (at \textbf{2170}) & 826 \\
\midrule
paper (digitised) & --- & ${\approx}5000$ & --- & \textbf{${\approx}1.3$} \\
\bottomrule
\end{tabular}
\caption{Landau--Zener at $v_q^D/\vth = 1$, $\Npc = 2.7\times10^3$, $\wpt = 2\times10^4$.
The paper's value is read off its figure: the red PIC curve moves 17 pixels over the
full $x$-range at ${\sim}153$ px per decade, i.e. a factor $10^{0.11} = 1.3$, which is
its ``the kinetic energy of the electrons is not even changed by a factor of 2''.}
\end{table}

\textbf{We measure ${\sim}10^3$ where the paper reports ${\lesssim}2$, and the cause
is not yet identified.} Four candidates have been tested and excluded:

\begin{enumerate}[nosep]
  \item \textbf{The loader} (\S\ref{sec:quietrelax}). $1003 \rightarrow 962$, $-4\%$.
        Predicted in advance from \code{fast}, where the loader moved the peak field
        energy $4.8\times$ but $\Delta KE_e$ only $7\%$.
  \item \textbf{The sweep convention} (\S\ref{sec:sweep}). $962 \rightarrow 826$,
        $-14\%$. The convention demonstrably took effect --- the peak moved from
        $\wpt = 4269$ to $2170$, halving as a doubled sweep rate requires --- so this
        is a real null, not a failed switch.
  \item \textbf{Coherent versus random energy.} $T^{00}$ counts the uniform $k=0$
        quiver motion, which is sloshing rather than heating, and at
        $v_q^D/\vth = 1$ the forced oscillation alone is $2.78\,\vth$. Decomposing
        into $\tfrac12\langle u\rangle^2$ and
        $\tfrac12(\langle u^2\rangle - \langle u\rangle^2)$ gives a coherent fraction
        never above $0.5\%$; the random ratio equals the total ratio to three digits.
        The heating is genuinely thermal.
  \item \textbf{The drive implementation.} \code{e\_ext} matches
        $A_0\cos(\omega_0 t + \tfrac12\dot\omega t^2)$ to $9.5\times10^{-15}$ once the
        \code{CustomStat} $t-\Delta t$ offset is applied, and does not match the naive
        form at all. Amplitude, sweep rate and waveform are all as intended.
\end{enumerate}

Against this, both resonance cases now agree with the paper (\S\ref{sec:resonance}),
so the machinery is sound and the disagreement is specific to the swept drive. One
diagnostic: by the end of the run $u_{\rm rms}\approx1.2c$, so our electrons leave the
non-relativistic regime the paper works in --- a symptom of the excess energy rather
than an explanation for it.

\subsection{Particle counts, and why the low-amplitude case stays gated}

Before resonance the swept drive is a \emph{bounded} forced oscillation of amplitude
$A_0/(1-\omega^2/\omega_p^2) = 2.78A_0$ at $\omega = 0.8\omega_p$, not a growing one,
so there is no time at which the signal outgrows the noise. The paper requires
$\bar t_{\rm noise} = L/(A_0\sqrt{3N_pN_x/2}) \le 0.02$, giving $\Npc = 2.7\times10^3$
and $2.7\times10^5$ for the two amplitudes. Verified for the run used here:
$\Npc = 2700$ gives $\bar t_{\rm noise} = 0.0199$.

The $v_q^D/\vth = 0.1$ case is written (\code{run\_lz0p1.sh}, $\Npc = 2.7\times10^5$,
$5.3\times10^8$ particles ${\approx}31$\,GB, three A100s, ${\sim}81$\,h chained) and
deliberately \textbf{not} run. It would cost ${\sim}243$ GPU-hours to produce a second
number against the same unresolved normalisation.

\section{Superseded results}
\label{sec:corrections}

Figures from earlier drafts of this document, and from comments still present in some
older scripts, that are wrong. Listed so they are recognisable if encountered --- not
as a history.

\begin{center}\small
\begin{tabular}{@{}p{5.0cm}p{9.4cm}@{}}
\toprule
Superseded claim & Correct statement \\
\midrule
``The quiet start is validated: $\delta n/n = 0$ exactly, noise floor
$1.31\times10^{-17}$ against $3.03\times10^{-6}$ random.''
& Those gates measured the wrong property. The paper's initial condition requires
  $\epsilon_{\rm noise}(0)=0$ \emph{and} relaxation to the Poisson floor within a few
  $\omega_p^{-1}$ (its Fig.~A00). The legacy loader satisfied only the first: keying
  velocity on the class makes every class a rigid comb, so the lattice never
  decohered and driven runs seeded daughter modes from round-off
  (\S\ref{sec:quietrelax}). ``Noise floor $10^{-17}$'' was the defect, reported as
  an achievement. \\
\addlinespace
``The failed first design gave $\delta n/n = 6.3\times10^{-3}$ after one step.''
& That is $1/\sqrt{\Npc}$ --- physical shot noise, and the paper's behaviour. It was
  engineered away. \\
\addlinespace
``Both resonance cases overshoot the thermal energy in the field energy (43.6 versus
2.1).''
& An artefact of the legacy loader. With the sampling loader the fast case peaks at
  $9.14\,\epsilon_{\rm th}$ and the slow case at $0.53$, never crossing the thermal
  line --- which is what the paper describes (\S\ref{sec:resonance}). \\
\addlinespace
``The fast case should peak near $\wpt \approx 110$, as seen in earlier random-start
scans.''
& $110$ is the analytic estimate $t_{\rm sat} + 1/\omega_{pi}$, not a measurement,
  and was mislabelled as observed. The random-start scan (\code{run34} \code{r3e-2})
  peaks at $\wpt = 241$. The sampling loader gives $238$ --- agreement with its
  control to $1\%$. \\
\addlinespace
``The paper's $\omega_p/H\approx7\times10^4$ shows it used the integrated sweep
phase.''
& Its figure places the crossing near $\wpt = 5000$, which is the naive convention;
  the quoted $\omega_p/H$ is then too large by two (\S\ref{sec:sweep}). This argument
  was used to cancel the naive-sweep test before it ran. \\
\addlinespace
``The shot-noise floor is $650\times$ below Eq.~B4.''
& A units error: $N_p$ is the \emph{total} macroparticle count, not per cell.
  Correctly evaluated, Eq.~B4 agrees with measurement at a flat $3.1\times$ and is
  the right tool for sizing (\S\ref{sec:sizing}). \\
\addlinespace
``Per-step cost is stable.''
& It roughly doubles at saturation (\S\ref{sec:slowdown}). A benchmark confined to
  the laminar phase measures the wrong number; the \code{slow} re-run took 43\,h
  against a $25$\,h estimate for the same reason. \\
\bottomrule
\end{tabular}
\end{center}

\noindent Two recurring failure modes are worth naming. The first is extrapolating a
measurement outside the regime where it was taken: the cost model, calibrated at
$\Npc \le 8000$, predicted 8 minutes for a run that took 41 and 97 seconds for one
that exceeded 40. The second is \textbf{validating against a property that is easy to
measure rather than the one that matters} --- the quiet-start gates checked that the
lattice was perfect, when the requirement was that it stop being perfect promptly.
Both survive a long time because the wrong quantity looks excellent.

\section{Status and remaining work}
\label{sec:status}

\paragraph{Complete.} The quiet-start loader and its gate against the paper's
Fig.~A00; both resonance cases of Fig.~2 with the corrected loader, \emph{both now
agreeing with the paper}; the Landau--Zener case at $v_q^D/\vth = 1$ under both sweep
conventions; MPI decomposition independence; the full performance
characterisation; and the reduction of every completed run to \code{.npz}
(\code{extract\_paper.py}, ${\sim}4$--21 min on one \code{vm-small} node, 50\,MB).

\begin{center}\small
\begin{tabular}{@{}llrl@{}}
\toprule
Run & tag & wall & outcome \\
\midrule
loader gate            & \code{gate\_*}   & 20\,min  & at the Poisson floor by $\wpt=1$ \\
resonant fast          & \code{fast\_s}   & 31\,min  & $kT_e \times 75.5$, paper says $O(10^2)$ \\
resonant slow          & \code{slow\_s}   & 43\,h    & peak $0.53\,\epsilon_{\rm th}$, never crosses thermal \\
LZ, integrated sweep   & \code{lz1p0\_s}  & 2:46     & $kT_e \times 962$ \\
LZ, naive sweep        & \code{lz1p0n}    & 3:03     & $kT_e \times 826$ \\
\bottomrule
\end{tabular}
\end{center}

\paragraph{Blocking.} The Landau--Zener heating discrepancy: ${\sim}10^3$ against the
paper's ${\approx}1.3$ (\S\ref{sec:lz}). Four candidates tested and excluded ---
loader, sweep convention, coherent/random decomposition, drive implementation. Since
both resonance cases now reproduce, the fault is specific to the swept drive rather
than to the code generally. The $v_q^D/\vth = 0.1$ case stays gated: ${\sim}243$
GPU-hours for a second number against the same unresolved normalisation.

\paragraph{Open.}
\begin{itemize}[nosep]
  \item \textbf{The \code{heat} gate predates the loader fix.} The numerical-heating
        bound over $10^6$ steps ($\Delta T_e = -0.016\%$) was measured with the legacy
        loader, whose lattice cannot heat because it cannot decohere. It needs
        re-running (${\sim}2$\,h) before the Landau--Zener results rest on it.
  \item \textbf{Only the Landau--Zener figure has been digitised.} The resonance
        comparison still rests on the paper's text. The published Fig.~2 would give
        peak $\epsilon_E/\epsilon_{\rm th}$ and its timing directly, which are the
        sharpest discriminators available.
  \item \textbf{Long-term stability} ($\wpt \to 8\times10^4$), the paper's strongest
        claim, is untested; at $\Npc = 2\times10^5$ it is roughly eight times the slow
        case, which already took 43\,h on three A100s.
  \item \textbf{2D cross-check} not run; 1+1D sufficiency is assumed on the paper's
        word.
  \item \textbf{Collisions.} \textsc{Entity}, like \textsc{SHARP}, is collisionless.
        ``Persists indefinitely'' should be read as ``on collisionless timescales''.
\end{itemize}

\end{document}
